\documentclass[aspectratio=169, 10pt]{beamer}

\usepackage[utf8]{inputenc}
\usepackage[T5]{fontenc}
\usepackage[vietnamese]{babel}
\usepackage{amsmath, amssymb, amsfonts}
\usepackage{array, tabularx, booktabs}
\usepackage{xcolor}
\usepackage{tikz}
\usepackage[skins,breakable]{tcolorbox}

% Thiết lập giao diện Beamer
\setbeamertemplate{navigation symbols}{}
\setbeamercolor{background canvas}{bg=white}
\setbeamertemplate{itemize item}{\color{cozyWarm}$\blacktriangleright$}
\setbeamertemplate{itemize subitem}{\color{cozySlate}$\bullet$}

% Bảng màu Cozy
\definecolor{cozySlate}{RGB}{44, 62, 80}       % #2C3E50 (Tiêu đề, cấu trúc)
\definecolor{cozyBg}{RGB}{250, 248, 245}       % #FAF8F5 (Nền hộp ấm)
\definecolor{cozyWarm}{RGB}{211, 84, 0}        % #D35400 (Cam đất nhấn mạnh)
\definecolor{cozyAmber}{RGB}{230, 126, 34}      % #E67E22 (Amber phụ)
\definecolor{cozyTrue}{RGB}{46, 125, 50}       % #2E7D32 (Đúng / 1)
\definecolor{cozyFalse}{RGB}{192, 57, 43}      % #C0392B (Sai / 0)
\definecolor{cozyMuted}{RGB}{127, 140, 141}    % #7F8C8D (Ghi chú)
\definecolor{cozyLightBlue}{RGB}{235, 243, 250}
\definecolor{cozyLightGreen}{RGB}{235, 247, 238}
\definecolor{cozyLightRed}{RGB}{253, 237, 236}
\definecolor{cozyLightAmber}{RGB}{254, 245, 231}

\setbeamercolor{frametitle}{fg=cozySlate}
\setbeamerfont{frametitle}{series=\bfseries, size=\large}

% Footline 3 khối màu Cozy
\setbeamercolor{footleftbox}{bg=cozySlate, fg=white}
\setbeamercolor{footmidbox}{bg=cozyWarm, fg=white}
\setbeamercolor{footrightbox}{bg=cozySlate!85!black, fg=white}

\setbeamertemplate{footline}{%
  \leavevmode%
  \hbox{%
    \begin{beamercolorbox}[wd=.333333\paperwidth,ht=2.3ex,dp=1ex,left]{footleftbox}%
      \hspace*{2ex}\scriptsize\bfseries Cấu Trúc Rời Rạc -- UIT/CITD%
    \end{beamercolorbox}%
    \begin{beamercolorbox}[wd=.333333\paperwidth,ht=2.3ex,dp=1ex,center]{footmidbox}%
      \scriptsize\bfseries Chương 1: Logic Mệnh Đề \& Suy Luận%
    \end{beamercolorbox}%
    \begin{beamercolorbox}[wd=.333334\paperwidth,ht=2.3ex,dp=1ex,right]{footrightbox}%
      \scriptsize\bfseries \insertframenumber{} / \inserttotalframenumber\hspace*{2ex}%
    \end{beamercolorbox}%
  }%
  \vskip0pt%
}

% 5 Môi trường tcolorbox chuyên biệt
\newtcolorbox{cozyDef}[1][]{
  enhanced, colback=white, colframe=cozySlate, coltitle=white,
  fonttitle=\bfseries\small, title={#1}, arc=2mm, boxrule=0.8pt,
  left=2mm, right=2mm, top=1.5mm, bottom=1.5mm
}

\newtcolorbox{cozyInsight}[1][]{
  enhanced, colback=cozyLightGreen, colframe=cozyTrue, coltitle=white,
  fonttitle=\bfseries\small, title={#1}, arc=2mm, boxrule=0.8pt,
  left=2mm, right=2mm, top=1.5mm, bottom=1.5mm
}

\newtcolorbox{cozyTrap}[1][]{
  enhanced, colback=cozyLightRed, colframe=cozyFalse, coltitle=white,
  fonttitle=\bfseries\small, title={#1}, arc=2mm, boxrule=0.8pt,
  left=2mm, right=2mm, top=1.5mm, bottom=1.5mm
}

\newtcolorbox{cozyMethod}[1][]{
  enhanced, colback=cozyLightAmber, colframe=cozyWarm, coltitle=white,
  fonttitle=\bfseries\small, title={#1}, arc=2mm, boxrule=0.8pt,
  left=2mm, right=2mm, top=1.5mm, bottom=1.5mm
}

\newtcolorbox{cozyExample}[1][]{
  enhanced, colback=cozyBg, colframe=cozyMuted, coltitle=white,
  fonttitle=\bfseries\small, title={#1}, arc=2mm, boxrule=0.6pt,
  left=2mm, right=2mm, top=1.5mm, bottom=1.5mm
}

\begin{document}

\begin{frame}[plain]
  \centering
  \vspace*{1.2cm}
  {\Huge\textbf{\color{cozySlate}CẤU TRÚC RỜI RẠC}}
  \par\vspace{0.3cm}
  {\color{cozyWarm}\rule{0.65\textwidth}{1.8pt}}
  \par\vspace{0.4cm}
  {\LARGE\textbf{\color{cozyWarm}Chương 1: Logic Mệnh Đề \& Phương Pháp Suy Luận}}
  \par\vspace{0.6cm}
  {\normalsize\textbf{\color{cozySlate}Tài liệu học tập trực quan -- Tương tác -- Dễ hiểu}}
  \par\vspace{0.3cm}
  {\scriptsize\color{cozyMuted}Khoa Công nghệ Thông tin -- Đại học Quốc gia TP.HCM}
\end{frame}

\begin{frame}{Mục tiêu \& Lộ trình Học tập Chương 1}
  \begin{columns}[t, onlytextwidth]
    \begin{column}{0.48\textwidth}
      \begin{cozyDef}[1. Mệnh đề \& Phân loại]
        Hiểu bản chất nhị phân của chân lý ($1/0$), phân biệt mệnh đề sơ cấp vs phức hợp.
      \end{cozyDef}
      \vspace{0.2cm}
      \begin{cozyInsight}[2. Các Toán tử Logic]
        Làm chủ 6 phép toán logic, giải mã phép kéo theo $p \to q$ qua mô hình lời hứa.
      \end{cozyInsight}
    \end{column}
    \begin{column}{0.48\textwidth}
      \begin{cozyMethod}[3. Tương đương Logic]
        Hệ thống 11 luật theo 3 cụm nhớ lâu, quy trình biến đổi đại số mệnh đề 3 bước.
      \end{cozyMethod}
      \vspace{0.2cm}
      \begin{cozyTrap}[4. Quy tắc Suy luận \& Thực tế]
        8 quy tắc suy luận, phân biệt ngụy biện bài thi, giải mã bài toán dã ngoại 8 bước.
      \end{cozyTrap}
    \end{column}
  \end{columns}
\end{frame}

% ==========================================
% PHẦN 1: KHÁI NIỆM MỆNH ĐỀ & PHÂN LOẠI
% ==========================================

\begin{frame}[plain]
  \centering
  \vspace*{2.0cm}
  {\LARGE\textbf{\color{cozySlate}1. Khái Niệm Mệnh Đề \& Phân Loại}}
  \par\vspace{0.4cm}
  {\color{cozyWarm}\rule{0.6\textwidth}{1.5pt}}
  \par\vspace{0.3cm}
  {\normalsize\color{cozyMuted}Nền tảng của logic hình thức và toán học rời rạc}
\end{frame}

\begin{frame}{Mệnh Đề là Gì? (Bản chất Nhị phân)}
  \begin{cozyDef}[Định nghĩa Mệnh đề (Proposition)]
    Một \textbf{mệnh đề} là một câu khẳng định một sự kiện khách quan, có giá trị chân lý hoặc \textbf{Đúng} ($1$) hoặc \textbf{Sai} ($0$), tuyệt đối không thể vừa đúng vừa sai hoặc không xác định.
  \end{cozyDef}
  \vspace{0.05cm}
  \begin{columns}[t, onlytextwidth]
    \begin{column}{0.48\textwidth}
      \begin{cozyInsight}[Ví dụ về Mệnh Đề]
        \small
        \begin{itemize}
          \setlength{\itemsep}{1pt}
          \item ``Hà Nội là thủ đô của Việt Nam.''\\
          $\to$ \textbf{\textcolor{cozyTrue}{Đúng (1)}}
          \item ``$0 > 2$''\\
          $\to$ \textbf{\textcolor{cozyFalse}{Sai (0)}}
          \item ``$3$ là một số nguyên tố.''\\
          $\to$ \textbf{\textcolor{cozyTrue}{Đúng (1)}}
        \end{itemize}
      \end{cozyInsight}
    \end{column}
    \begin{column}{0.48\textwidth}
      \begin{cozyTrap}[KHÔNG Phải Mệnh Đề]
        \small
        \begin{itemize}
          \setlength{\itemsep}{1pt}
          \item ``Một ngày đẹp trời làm sao!''\\
          \textcolor{cozyMuted}{(Câu cảm thán)}
          \item ``Thầy Hợi rất đẹp trai!''\\
          \textcolor{cozyMuted}{(Đánh giá chủ quan)}
          \item ``$3 + x = 2$''\\
          \textcolor{cozyMuted}{(Vị từ chứa biến, chưa biết $x$)}
          \item ``Câu này là sai.''\\
          \textcolor{cozyMuted}{(Nghịch lý tự quy chiếu)}
        \end{itemize}
      \end{cozyTrap}
    \end{column}
  \end{columns}
\end{frame}

\begin{frame}{Ký Hiệu, Chân Trị \& Góc Nhìn Khoa Học Máy Tính}
  \begin{columns}[t, onlytextwidth]
    \begin{column}{0.48\textwidth}
      \begin{cozyDef}[Ký hiệu \& Chân trị (Truth Value)]
        \small
        \begin{itemize}
          \item \textbf{Ký hiệu:} $p, q, r, s, \dots$
          \item \textbf{Chân trị:} Mỗi mệnh đề nhận duy nhất một trong hai giá trị:
          \begin{itemize}
            \item \textbf{Đúng} (True): Ký hiệu \textbf{Đ} hoặc $\mathbf{1}$.
            \item \textbf{Sai} (False): Ký hiệu \textbf{S} hoặc $\mathbf{0}$.
          \end{itemize}
          \item Tuyệt đối không có trạng thái thứ ba.
        \end{itemize}
      \end{cozyDef}
    \end{column}
    \begin{column}{0.48\textwidth}
      \begin{cozyInsight}[Góc nhìn Khoa học Máy tính]
        \small
        \begin{itemize}
          \item Trong máy tính, chân lý nhị phân chính là \textbf{Bit ($1 / 0$)}, là mức điện thế Cao/Thấp trong vi mạch bán dẫn.
          \item Trong lập trình, đây chính là kiểu dữ liệu nguyên thủy \textbf{\texttt{bool}} (\texttt{True} / \texttt{False}), quyết định mọi luồng rẽ nhánh \texttt{if / else}, vòng lặp và logic điều khiển trong phần mềm.
        \end{itemize}
      \end{cozyInsight}
    \end{column}
  \end{columns}
\end{frame}

\begin{frame}{Phân Loại Mệnh Đề: Sơ Cấp vs Phức Hợp}
  \begin{cozyDef}[1. Mệnh đề Sơ cấp (Nguyên tử -- Atomic Proposition)]
    \small
    Chỉ chứa duy nhất một sự kiện, không thể chia nhỏ hơn thành các mệnh đề khác.
    \begin{itemize}
      \item \textit{Ví dụ:} ``$3$ là số nguyên tố'', ``$2$ là số chẵn''.
    \end{itemize}
  \end{cozyDef}
  \vspace{0.15cm}
  \begin{cozyMethod}[2. Mệnh đề Phức hợp (Compound Proposition)]
    \small
    Được hình thành bằng cách kết hợp các mệnh đề sơ cấp lại với nhau thông qua các \textbf{toán tử logic} (liên từ: \textit{và, hoặc, phủ định, nếu\dots thì, nếu và chỉ nếu}).
    \begin{itemize}
      \item \textit{Ví dụ 1:} ``$2$ là số nguyên tố \textbf{\color{cozyWarm}VÀ} là một số chẵn.''
      \item \textit{Ví dụ 2:} ``Thầy Hợi đang dạy học \textbf{\color{cozyWarm}HOẶC} đang xem phim.''
    \end{itemize}
    \vspace{0.1cm}
    \par
    \textbf{\color{cozySlate}Lưu ý sư phạm:} Trong toán học, mệnh đề phức hợp được mô hình hóa như một hàm chân trị nhận các biến mệnh đề đầu vào.
  \end{cozyMethod}
\end{frame}

% ==========================================
% PHẦN 2: 6 TOÁN TỬ LOGIC CƠ BẢN
% ==========================================

\begin{frame}[plain]
  \centering
  \vspace*{2.0cm}
  {\LARGE\textbf{\color{cozySlate}2. Các Toán Tử Logic Cơ Bản}}
  \par\vspace{0.4cm}
  {\color{cozyWarm}\rule{0.6\textwidth}{1.5pt}}
  \par\vspace{0.3cm}
  {\normalsize\color{cozyMuted}Ngôn ngữ hình thức để xây dựng các mệnh đề phức hợp}
\end{frame}

\begin{frame}{1. Phép Phủ Định (NOT -- $\neg p$ hoặc $\bar{p}$)}
  \begin{columns}[t, onlytextwidth]
    \begin{column}{0.44\textwidth}
      \begin{cozyDef}[Định nghĩa \& Bảng Chân Lý]
        \small
        Phủ định của mệnh đề $p$ là mệnh đề \textbf{``Không phải trường hợp $p$''}, ký hiệu là $\neg p$ hoặc $\bar{p}$.\\
        \vspace{0.05cm}
        Giá trị chân lý của $\bar{p}$ luôn \textbf{ngược lại} với $p$: đúng khi $p$ sai, và sai khi $p$ đúng.
        \vspace{0.2cm}
        \begin{center}
          \renewcommand{\arraystretch}{1.15}
          \begin{tabular}{cc}
            \toprule
            \textbf{Mệnh đề ($p$)} & \textbf{Phủ định ($\bar{p}$)} \\
            \midrule
            $0$ & \textbf{\textcolor{cozyTrue}{1}} \\
            $1$ & \textbf{\textcolor{cozyFalse}{0}} \\
            \bottomrule
          \end{tabular}
        \end{center}
        \vspace{0.1cm}
        \par\footnotesize\centering\color{cozyMuted}
        Đảo bit: $0 \leftrightarrow 1$ (Cổng NOT logic)
      \end{cozyDef}
    \end{column}
    \begin{column}{0.53\textwidth}
      \begin{cozyInsight}[Ví dụ Minh Họa từ Thực Tế]
        \small
        \begin{itemize}
          \setlength{\itemsep}{4pt}
          \item $p$: ``Paris là thủ đô của Pháp.'' (\textbf{\textcolor{cozyTrue}{1}})\\
          $\implies \bar{p}$: ``Paris \textbf{không phải} là thủ đô của Pháp.'' (\textbf{\textcolor{cozyFalse}{0}})
          \item $q$: ``$9$ là một số chính phương.'' (\textbf{\textcolor{cozyTrue}{1}})\\
          $\implies \bar{q}$: ``$9$ \textbf{không phải} là một số chính phương.'' (\textbf{\textcolor{cozyFalse}{0}})
          \item $r$: ``$0 > 2$'' (\textbf{\textcolor{cozyFalse}{0}})\\
          $\implies \bar{r}$: ``$0 \le 2$'' (\textbf{\textcolor{cozyTrue}{1}})
        \end{itemize}
      \end{cozyInsight}
    \end{column}
  \end{columns}
\end{frame}

\begin{frame}{2. Phép Tuyển Thường (OR -- $p \lor q$)}
  \begin{columns}[t, onlytextwidth]
    \begin{column}{0.44\textwidth}
      \begin{cozyDef}[Định nghĩa \& Bảng Chân Lý]
        \small
        Phép tuyển của $p$ và $q$ là mệnh đề \textbf{``$p$ hoặc $q$''}, ký hiệu là $p \lor q$.\\
        \vspace{0.05cm}
        $p \lor q$ \textbf{\textcolor{cozyTrue}{Đúng}} khi có \textbf{ít nhất một} mệnh đề thành phần đúng; chỉ \textbf{\textcolor{cozyFalse}{Sai}} khi cả $p$ và $q$ đều sai.
        \vspace{0.15cm}
        \begin{center}
          \renewcommand{\arraystretch}{1.1}
          \begin{tabular}{cc|c}
            \toprule
            $p$ & $q$ & $p \lor q$ \\
            \midrule
            $0$ & $0$ & \textbf{\textcolor{cozyFalse}{0}} \\
            $0$ & $1$ & \textbf{\textcolor{cozyTrue}{1}} \\
            $1$ & $0$ & \textbf{\textcolor{cozyTrue}{1}} \\
            $1$ & $1$ & \textbf{\textcolor{cozyTrue}{1}} \\
            \bottomrule
          \end{tabular}
        \end{center}
        \vspace{0.05cm}
        \par\footnotesize\centering\color{cozyMuted}Chỉ cần 1 đầu vào đúng $\implies$ Đúng
      \end{cozyDef}
    \end{column}
    \begin{column}{0.53\textwidth}
      \begin{cozyInsight}[Bản chất Inclusive OR \& Ví dụ]
        \small
        \begin{itemize}
          \setlength{\itemsep}{3pt}
          \item \textbf{Inclusive OR (Tuyển bao hàm):} Trong toán học, ``hoặc'' cho phép \textbf{cả hai cùng xảy ra}.
          \item \textit{Ví dụ xét học bổng:} ``Sinh viên có GPA $\ge 8.0$ \textbf{HOẶC} có bài báo nghiên cứu khoa học.''\\
          $\to$ Đạt \textbf{cả hai} điều kiện thì càng xứng đáng!
          \item \textit{Ví dụ slide gốc:} ``$3$ là số nguyên tố \textbf{hoặc} $3$ là số chẵn.''\\
          $\to$ Vì $3$ là số nguyên tố (Đúng) $\implies$ Toàn bộ mệnh đề \textbf{\textcolor{cozyTrue}{Đúng (1)}}.
        \end{itemize}
      \end{cozyInsight}
    \end{column}
  \end{columns}
\end{frame}

\begin{frame}{3. Phép Hội (AND -- $p \land q$)}
  \begin{columns}[t, onlytextwidth]
    \begin{column}{0.44\textwidth}
      \begin{cozyDef}[Định nghĩa \& Bảng Chân Lý]
        \small
        Phép hội của $p$ và $q$ là mệnh đề \textbf{``$p$ và $q$''}, ký hiệu là $p \land q$.\\
        \vspace{0.05cm}
        $p \land q$ chỉ \textbf{\textcolor{cozyTrue}{Đúng}} khi \textbf{cả hai} mệnh đề đều đúng; chỉ cần một mệnh đề sai thì kết quả sẽ \textbf{\textcolor{cozyFalse}{Sai}}.
        \vspace{0.15cm}
        \begin{center}
          \renewcommand{\arraystretch}{1.1}
          \begin{tabular}{cc|c}
            \toprule
            $p$ & $q$ & $p \land q$ \\
            \midrule
            $0$ & $0$ & \textbf{\textcolor{cozyFalse}{0}} \\
            $0$ & $1$ & \textbf{\textcolor{cozyFalse}{0}} \\
            $1$ & $0$ & \textbf{\textcolor{cozyFalse}{0}} \\
            $1$ & $1$ & \textbf{\textcolor{cozyTrue}{1}} \\
            \bottomrule
          \end{tabular}
        \end{center}
        \vspace{0.05cm}
        \par\footnotesize\centering\color{cozyMuted}Cả hai cùng đúng $\implies$ Mới đúng
      \end{cozyDef}
    \end{column}
    \begin{column}{0.53\textwidth}
      \begin{cozyMethod}[Đối Chiếu Trực Quan: $\lor$ (OR) vs $\land$ (AND)]
        \small
        \begin{itemize}
          \setlength{\itemsep}{3pt}
          \item \textbf{Phép $\lor$ rất ``dễ tính'':} Chỉ cần ít nhất $1$ nhánh đúng là đúng (mạch điện \textit{song song}).
          \item \textbf{Phép $\land$ rất ``khắt khe'':} Bắt buộc cả hai nhánh phải đúng (mạch điện \textit{nối tiếp}).
          \item \textit{Ví dụ slide gốc:} ``Trời đã mưa \textbf{VÀ} hôm đó là thứ Sáu.''\\
          $\to$ Chỉ đúng khi đồng thời cả 2 sự kiện xảy ra; sai 1 sự kiện $\implies$ phát biểu sai.
        \end{itemize}
      \end{cozyMethod}
    \end{column}
  \end{columns}
\end{frame}

\begin{frame}{4. Phép Tuyển Loại (XOR -- $p \oplus q$)}
  \begin{columns}[t, onlytextwidth]
    \begin{column}{0.44\textwidth}
      \begin{cozyDef}[Định nghĩa \& Bảng Chân Lý]
        \small
        Phép tuyển loại của $p$ và $q$ là mệnh đề \textbf{``$p$ hoặc $q$ nhưng không cả hai''}, ký hiệu $p \oplus q$.\\
        \vspace{0.05cm}
        $p \oplus q$ \textbf{\textcolor{cozyTrue}{Đúng}} khi $p, q$ có \textbf{chân trị khác nhau} (chỉ đúng $1$ trong $2$); \textbf{\textcolor{cozyFalse}{Sai}} khi cùng chân trị.
        \vspace{0.15cm}
        \begin{center}
          \renewcommand{\arraystretch}{1.1}
          \begin{tabular}{cc|c}
            \toprule
            $p$ & $q$ & $p \oplus q$ \\
            \midrule
            $0$ & $0$ & \textbf{\textcolor{cozyFalse}{0}} \\
            $0$ & $1$ & \textbf{\textcolor{cozyTrue}{1}} \\
            $1$ & $0$ & \textbf{\textcolor{cozyTrue}{1}} \\
            $1$ & $1$ & \textbf{\textcolor{cozyFalse}{0}} \;\textcolor{cozyWarm}{$\leftarrow$ Khác $\lor$} \\
            \bottomrule
          \end{tabular}
        \end{center}
        \vspace{0.05cm}
        \par\footnotesize\centering\color{cozyMuted}Cộng nhị phân không nhớ: $1 \oplus 1 = 0$
      \end{cozyDef}
    \end{column}
    \begin{column}{0.53\textwidth}
      \begin{cozyTrap}[Phân biệt Cốt lõi: Tuyển thường $\lor$ vs Tuyển loại $\oplus$]
        \small
        \begin{itemize}
          \setlength{\itemsep}{3pt}
          \item \textbf{$\lor$ (Inclusive OR):} Cho phép chọn cả hai. Ví dụ: ``Biết tiếng Anh hoặc tiếng Nhật''.
          \item \textbf{$\oplus$ (Exclusive OR):} Bắt buộc chỉ chọn duy nhất 1 trong 2:
          \begin{itemize}
            \setlength{\itemsep}{2pt}
            \item \textit{Thực đơn cơm trưa:} ``Chọn canh chua HOẶC canh rau'' (chỉ được chọn đúng 1 món).
            \item \textit{Ví dụ slide gốc:} ``Chọn súp miễn phí hoặc salad miễn phí kèm bữa ăn.'' Nếu cố lấy cả hai $\implies$ Phạm quy (\textbf{\textcolor{cozyFalse}{Sai}})!
          \end{itemize}
        \end{itemize}
      \end{cozyTrap}
    \end{column}
  \end{columns}
\end{frame}

\begin{frame}{5. Phép Kéo Theo (Conditional -- $p \to q$)}
  \begin{columns}[t, onlytextwidth]
    \begin{column}{0.44\textwidth}
      \begin{cozyDef}[Định nghĩa \& Bảng Chân Lý]
        \small
        Phép kéo theo của $p$ và $q$ là mệnh đề \textbf{``Nếu $p$ thì $q$''}, ký hiệu là $p \to q$.\\
        \vspace{0.05cm}
        Trong đó: $p$ là \textbf{Tiền đề} (Giả thiết), $q$ là \textbf{Kết luận} (Hệ quả).
        \vspace{0.15cm}
        \begin{center}
          \renewcommand{\arraystretch}{1.1}
          \begin{tabular}{cc|c}
            \toprule
            $p$ & $q$ & $p \to q$ \\
            \midrule
            $0$ & $0$ & \textbf{\textcolor{cozyTrue}{1}} \\
            $0$ & $1$ & \textbf{\textcolor{cozyTrue}{1}} \\
            $1$ & $0$ & \textbf{\textcolor{cozyFalse}{0}} \;\textcolor{cozyFalse}{$\leftarrow$ DUY NHẤT SAI} \\
            $1$ & $1$ & \textbf{\textcolor{cozyTrue}{1}} \\
            \bottomrule
          \end{tabular}
        \end{center}
        \vspace{0.05cm}
        \par\footnotesize\centering\color{cozyMuted}Tiền đề sai $\implies$ Luôn luôn Đúng
      \end{cozyDef}
    \end{column}
    \begin{column}{0.53\textwidth}
      \begin{cozyTrap}[Quy Tắc Vàng Ghi Nhớ $p \to q$]
        \small
        \begin{itemize}
          \setlength{\itemsep}{3pt}
          \item Mệnh đề kéo theo $p \to q$ chỉ \textbf{\textcolor{cozyFalse}{Sai}} trong \textbf{duy nhất một trường hợp}:\\
          \vspace{0.1cm}
          \centerline{\textbf{\textcolor{cozySlate}{Tiền đề Đúng ($p=1$)}} $\implies$ \textbf{\textcolor{cozyFalse}{Kết luận Sai ($q=0$)}}}
          \vspace{0.1cm}
          \item Trong 3 trường hợp còn lại, $p \to q$ luôn \textbf{\textcolor{cozyTrue}{Đúng ($1$)}}.
          \item \textit{Ví dụ slide gốc:} ``Nếu bạn học chăm chỉ ($p$), thì bạn sẽ qua môn ($q$).''
        \end{itemize}
      \end{cozyTrap}
    \end{column}
  \end{columns}
\end{frame}

\begin{frame}{Giải Mã Phép Kéo Theo: Vì sao Tiền đề Sai thì Luôn Đúng?}
  \begin{cozyInsight}[Mô hình Lời Hứa / Hợp Đồng Thưởng Điểm]
    \small
    Thầy giáo hứa với học sinh: \textit{``Nếu em đạt điểm 10 môn Toán ($p$), thầy sẽ tặng em máy tính ($q$).''}\\
    Trong trường hợp nào thì thầy giáo bị coi là \textbf{nói dối / thất hứa} ($p \to q$ có giá trị Sai)?
    \begin{itemize}
      \setlength{\itemsep}{2pt}
      \item \textbf{Trường hợp 1 ($p=1, q=1$):} Em được 10 điểm, thầy tặng máy $\implies$ Thầy \textbf{giữ lời} (\textbf{\textcolor{cozyTrue}{Đúng}}).
      \item \textbf{Trường hợp 2 ($p=1, q=0$):} Em được 10 điểm, thầy không tặng $\implies$ Thầy \textbf{thất hứa} (\textbf{\textcolor{cozyFalse}{Sai}}).
      \item \textbf{Trường hợp 3 \& 4 ($p=0$):} Em \textbf{KHÔNG} được 10 điểm! Lúc này dù thầy có tặng hay không tặng máy tính, thầy đều \textbf{KHÔNG THẤT HỨA} $\implies$ Lời hứa hiển nhiên \textbf{\textcolor{cozyTrue}{Đúng}}!
    \end{itemize}
  \end{cozyInsight}
  \vspace{0.05cm}
  \begin{cozyDef}[Ghi Nhớ Cốt Lõi: Chân Trị Chân Không (Vacuous Truth)]
    \small
    Khi tiền đề $p$ \textbf{Sai ($0$)}, mệnh đề kéo theo $p \to q$ tự động nhận giá trị \textbf{\textcolor{cozyTrue}{Đúng ($1$)}} bất kể $q$ Đúng hay Sai. Trong logic học, hiện tượng này được gọi là \textbf{Vacuous Truth} (Đúng chân không).
  \end{cozyDef}
\end{frame}

\begin{frame}{Cạm Bẫy Thuật Ngữ: Điều Kiện Cần vs Điều Kiện Đủ}
  \begin{columns}[t, onlytextwidth]
    \begin{column}{0.48\textwidth}
      \begin{cozyTrap}[Phân biệt Điều kiện Cần và Đủ]
        \small
        Trong mệnh đề kéo theo $p \to q$:
        \begin{itemize}
          \setlength{\itemsep}{3pt}
          \item \textbf{$p$ là điều kiện ĐỦ để có $q$}:\\
          Chỉ cần có $p$ là chắc chắn có $q$. Có $p$ là ``đủ'' để suy ra $q$.
          \item \textbf{$q$ là điều kiện CẦN để có $p$}:\\
          Muốn có $p$ thì bắt buộc phải có $q$. Nếu không có $q$ thì chắc chắn không thể có $p$ ($\neg q \to \neg p$).
        \end{itemize}
      \end{cozyTrap}
    \end{column}
    \begin{column}{0.48\textwidth}
      \begin{cozyMethod}[Công Thức Tương Đương Cốt Lõi: $p \to q \equiv \neg p \lor q$]
        \small
        \begin{itemize}
          \setlength{\itemsep}{3pt}
          \item \textbf{Ý nghĩa trực giác:}\\
          ``Để không vi phạm lời hứa, hoặc là không có tiền đề ràng buộc ($\neg p$), hoặc là đã thực hiện kết luận ($q$).''
          \item \textbf{Ví dụ đời sống:}\\
          ``Trời mưa ($p$) $\to$ Đường ướt ($q$)''.
          \begin{itemize}
            \item Trời mưa là điều kiện \textit{đủ} để làm đường ướt.
            \item Đường ướt là điều kiện \textit{cần} để khẳng định trời đã mưa.
          \end{itemize}
        \end{itemize}
      \end{cozyMethod}
    \end{column}
  \end{columns}
\end{frame}

\begin{frame}{6. Phép Tương Đương (Biconditional -- $p \leftrightarrow q$)}
  \begin{columns}[t, onlytextwidth]
    \begin{column}{0.44\textwidth}
      \begin{cozyDef}[Định nghĩa \& Bảng Chân Lý]
        \small
        Mệnh đề tương đương của $p$ và $q$ là \textbf{``$p$ nếu và chỉ nếu $q$''}, ký hiệu $p \leftrightarrow q$.\\
        \vspace{0.05cm}
        $p \leftrightarrow q$ \textbf{\textcolor{cozyTrue}{Đúng}} khi $p, q$ có \textbf{cùng chân trị} ($1 \leftrightarrow 1$ là $1$, $0 \leftrightarrow 0$ là $1$); \textbf{\textcolor{cozyFalse}{Sai}} khi khác chân trị.
        \vspace{0.15cm}
        \begin{center}
          \renewcommand{\arraystretch}{1.1}
          \begin{tabular}{cc|c}
            \toprule
            $p$ & $q$ & $p \leftrightarrow q$ \\
            \midrule
            $0$ & $0$ & \textbf{\textcolor{cozyTrue}{1}} \\
            $0$ & $1$ & \textbf{\textcolor{cozyFalse}{0}} \\
            $1$ & $0$ & \textbf{\textcolor{cozyFalse}{0}} \\
            $1$ & $1$ & \textbf{\textcolor{cozyTrue}{1}} \\
            \bottomrule
          \end{tabular}
        \end{center}
        \vspace{0.05cm}
        \par\footnotesize\centering\color{cozyMuted}Cùng chân trị $\implies 1$; Khác chân trị $\implies 0$
      \end{cozyDef}
    \end{column}
    \begin{column}{0.53\textwidth}
      \begin{cozyInsight}[Bản chất Kéo Theo 2 Chiều \& Ví dụ]
        \small
        \begin{itemize}
          \setlength{\itemsep}{3pt}
          \item \textbf{Bản chất là kéo theo 2 chiều:}
          $$p \leftrightarrow q \equiv (p \to q) \land (q \to p)$$
          $p$ vừa là điều kiện cần, vừa là điều kiện đủ của $q$.
          \item \textit{Ví dụ từ giáo trình:}\\
          ``Bạn có thể lên chuyến bay nếu và chỉ nếu bạn mua vé.''
          \begin{itemize}
            \setlength{\itemsep}{2pt}
            \item Có vé $\implies$ Được bay ($p \to q$).
            \item Được bay $\implies$ Đã có vé ($q \to p$).
          \end{itemize}
        \end{itemize}
      \end{cozyInsight}
    \end{column}
  \end{columns}
\end{frame}

\begin{frame}{Kiểm Tra Nhanh Phản Xạ Chân Trị (Phần 1)}
  \begin{cozyExample}[Bài toán: Cho $p = 1$ (Đúng) và $q = 0$ (Sai). Tính chân trị các biểu thức sau:]
    \small
    \begin{columns}[c, onlytextwidth]
      \begin{column}{0.48\textwidth}
        \centering\renewcommand{\arraystretch}{1.15}
        \begin{tabular}{rll}
          \toprule
          \textbf{STT} & \textbf{Biểu thức} & \textbf{Chân trị} \\
          \midrule
          \textbf{1.} & $\neg p = \neg 1$ & $= \mathbf{\textcolor{cozyFalse}{0}}$ \\
          \textbf{2.} & $p \lor q = 1 \lor 0$ & $= \mathbf{\textcolor{cozyTrue}{1}}$ \\
          \textbf{3.} & $p \land q = 1 \land 0$ & $= \mathbf{\textcolor{cozyFalse}{0}}$ \\
          \textbf{4.} & $p \oplus q = 1 \oplus 0$ & $= \mathbf{\textcolor{cozyTrue}{1}}$ \\
          \bottomrule
        \end{tabular}
      \end{column}
      \begin{column}{0.48\textwidth}
        \centering\renewcommand{\arraystretch}{1.15}
        \begin{tabular}{rll}
          \toprule
          \textbf{STT} & \textbf{Biểu thức} & \textbf{Chân trị} \\
          \midrule
          \textbf{5.} & $p \to q = 1 \to 0$ & $= \mathbf{\textcolor{cozyFalse}{0}}$ \\
          \textbf{6.} & $q \to p = 0 \to 1$ & $= \mathbf{\textcolor{cozyTrue}{1}}$ \\
          \textbf{7.} & $p \leftrightarrow q = 1 \leftrightarrow 0$ & $= \mathbf{\textcolor{cozyFalse}{0}}$ \\
          \multicolumn{3}{c}{\footnotesize\color{cozyMuted}Khác chân trị $\implies 0$} \\
          \bottomrule
        \end{tabular}
      \end{column}
    \end{columns}
  \end{cozyExample}
  \vspace{0.05cm}
  \begin{cozyTrap}[Điểm Nhấn Sư Phạm: Vì sao $p \to q \neq q \to p$?]
    \small
    \begin{itemize}
      \setlength{\itemsep}{2pt}
      \item $p \to q = 1 \to 0$: Tiền đề Đúng nhưng Kết luận Sai $\implies$ Thất hứa $\implies$ \textbf{\textcolor{cozyFalse}{Sai ($0$)}}.
      \item $q \to p = 0 \to 1$: Tiền đề Sai ($q=0$) $\implies$ Mệnh đề hiển nhiên \textbf{\textcolor{cozyTrue}{Đúng ($1$)}} (Vacuous Truth).
      \item \textbf{Quy tắc cốt lõi:} Phép kéo theo \textbf{không có tính giao hoán}: $p \to q \not\equiv q \to p$.
    \end{itemize}
  \end{cozyTrap}
\end{frame}

\begin{frame}{Kiểm Tra Nhanh Phản Xạ Chân Trị (Phần 2)}
  \begin{cozyExample}[Xét tính Đúng / Sai của các phát biểu kéo theo ngôn ngữ tự nhiên:]
    \small
    \begin{enumerate}
      \setlength{\itemsep}{4pt}
      \item \textbf{Phát biểu 1:} ``Nếu $0 > 2$ thì Paris là thủ đô của Pháp.''
      \begin{itemize}
        \setlength{\itemsep}{1pt}
        \item \textbf{Phân tích:} Tiền đề $p$: ``$0 > 2$'' (\textbf{\textcolor{cozyFalse}{Sai -- 0}}), Kết luận $q$: ``Paris là thủ đô Pháp'' (\textbf{\textcolor{cozyTrue}{Đúng -- 1}}).
        \item \textbf{Kết luận:} Vì tiền đề Sai $\implies$ Mệnh đề kéo theo \textbf{\textcolor{cozyTrue}{Đúng (1)}}! (Vacuous Truth).
      \end{itemize}

      \item \textbf{Phát biểu 2:} ``Nếu $2 + 2 = 4$ thì $1 + 1 = 3$.''
      \begin{itemize}
        \setlength{\itemsep}{1pt}
        \item \textbf{Phân tích:} Tiền đề $p$: ``$2 + 2 = 4$'' (\textbf{\textcolor{cozyTrue}{Đúng -- 1}}), Kết luận $q$: ``$1 + 1 = 3$'' (\textbf{\textcolor{cozyFalse}{Sai -- 0}}).
        \item \textbf{Kết luận:} Rơi vào trường hợp duy nhất $1 \to 0$ $\implies$ Mệnh đề \textbf{\textcolor{cozyFalse}{Sai (0)}}!
      \end{itemize}

      \item \textbf{Phát biểu 3:} ``Nếu Mặt trời mọc ở hướng Tây thì $1 + 1 = 2$.''
      \begin{itemize}
        \setlength{\itemsep}{1pt}
        \item \textbf{Phân tích:} Tiền đề $p$: ``Mặt trời mọc ở hướng Tây'' (\textbf{\textcolor{cozyFalse}{Sai -- 0}}).
        \item \textbf{Kết luận:} Tiền đề Sai $\implies$ Bất kể kết luận ra sao, mệnh đề hiển nhiên \textbf{\textcolor{cozyTrue}{Đúng (1)}}!
      \end{itemize}
    \end{enumerate}
  \end{cozyExample}
\end{frame}

\end{document}
